\documentclass[10pt,a4paper]{article}
\usepackage[utf8]{inputenc}
\usepackage{amsmath, amsthm, amssymb}
\usepackage{geometry}
\usepackage{enumitem}
\usepackage{bm}

\geometry{
  a4paper,
  left=20mm,
  right=20mm,
  top=25mm,
  bottom=25mm,
}

\begin{document}
\section*{Problem}
\begin{enumerate}[label=\fbox{\arabic*}]
  \item % Problem 1
        Find all real numbers $a$, for which the improper integral
        \[
          \iint_D \frac{y}{(1 + \sqrt{x^2+y^2} - x)^2(x^2+y^2)^a} dxdy
        \]
        converges. Here, $D = \{(x, y) \in \mathbb{R}^2 \mid x \ge 0, y \ge 0, x^2+y^2 \ge 1\}$.

  \item % Problem 2
        Let $V$ be an $\mathbb{R}$-vector space and let $v_1, v_2, v_3, v_4$ be linearly independent vectors in $V$. For $a \in \mathbb{R}$, we put:
        \begin{align*}
          w_1 &= v_1 - v_2 + v_3 \\
          w_2 &= v_1 + (a - 1)v_2 + v_3 - a v_4 \\
          w_3 &= 2v_1 - 2v_2 + (a + 2)v_3 \\
          w_4 &= v_1 - v_2 + v_3 + (a + 1)v_4.
        \end{align*}
        \begin{enumerate}[label=(\arabic*)]
          \item Determine all possible $a$ for which $w_1, w_2, w_3, w_4$ are linearly independent.
          \item For $a$ such that $w_1, w_2, w_3, w_4$ are linearly dependent, determine the dimension of the subspace of $V$, $\langle w_1, w_2, w_3, w_4 \rangle$ generated by $w_1, w_2, w_3, w_4$.
        \end{enumerate}

  \item % Problem 3
        Let $H$ be the subgroup
        \[
          H = \{(a, b, 0, 0) \mid a \in \mathbb{Z}/2\mathbb{Z}, b \in \mathbb{Z}/3\mathbb{Z}\}
        \]
        of the abelian group $G = \mathbb{Z}/2\mathbb{Z} \times \mathbb{Z}/3\mathbb{Z} \times \mathbb{Z}/3\mathbb{Z} \times \mathbb{Z}/4\mathbb{Z}$. Determine the number of subgroups $K$ of $G$ of order $12$ satisfying $G = K + H$, where
        \[
          K + H = \{ \alpha + \beta \mid \alpha \in K, \beta \in H \}.
        \]

  \item % Problem 4
        Show that the function
        \[
          f(x) = \int_0^\infty |x + t \sin t| e^{-t} dt
        \]
        defined on $\mathbb{R}$ admits a minimum value.

  \item % Problem 5
        Find the value of the improper integral
        \[
          \int_{-\infty}^\infty \frac{x^3 \sin x}{(x^2 + 1)^2} dx.
        \]

  \item % Problem 6
        Let $X$ be the subspace of the Euclidean space $\mathbb{R}^4$ defined by
        \[
          X = \{ (x, y, z, w) \in \mathbb{R}^4 \mid x^2 + y^2 + z^2 + w^2 = 1, \ x^3 + y^3 + z^3 + w^3 = 0 \}.
        \]
        Show that $X$ is a compact differentiable manifold.
\end{enumerate}

\newpage
\section*{Solution}
\begin{enumerate}[label=\fbox{\arabic*}]
  \item % Solution 1
        We transition to polar coordinates $x = r \cos \theta$ and $y = r \sin \theta$. The domain of integration $D$ corresponds to $r \ge 1$ and $0 \le \theta \le \frac{\pi}{2}$. The area element is $dxdy = r dr d\theta$.
        Substituting these into the integral, we obtain:
        \begin{align*}
          I &= \int_0^{\pi/2} \int_1^\infty \frac{r \sin \theta}{(1 + r - r \cos \theta)^2 (r^2)^a} r dr d\theta \\
          &= \int_0^{\pi/2} \int_1^\infty \frac{r^{2-2a} \sin \theta}{(1 + r(1 - \cos \theta))^2} dr d\theta.
        \end{align*}
        Since the integrand is non-negative, we can safely apply Tonelli's theorem and perform the integration in either order or use substitution directly. 
        Let us substitute $c = 1 - \cos \theta$. As $\theta$ goes from $0$ to $\frac{\pi}{2}$, $c$ goes from $0$ to $1$. The differential is $dc = \sin \theta d\theta$. The integral transforms seamlessly into:
        \[
          I = \int_0^1 \int_1^\infty \frac{r^{2-2a}}{(1 + rc)^2} dr dc = \int_1^\infty r^{2-2a} \left( \int_0^1 \frac{1}{(1 + rc)^2} dc \right) dr.
        \]
        We evaluate the inner integral with respect to $c$:
        \[
          \int_0^1 \frac{1}{(1 + rc)^2} dc = \left[ -\frac{1}{r(1+rc)} \right]_{c=0}^{c=1} = -\frac{1}{r(1+r)} - \left(-\frac{1}{r}\right) = \frac{1}{r} - \frac{1}{r(1+r)} = \frac{1}{1+r}.
        \]
        Substituting this back into the outer integral:
        \[
          I = \int_1^\infty r^{2-2a} \left( \frac{1}{1+r} \right) dr = \int_1^\infty \frac{r^{2-2a}}{1+r} dr.
        \]
        Since the domain of integration is $[1, \infty)$, there are no singularities at $r = 0$. The integral converges if and only if it converges at the upper limit $r \to \infty$.
        As $r \to \infty$, the integrand behaves asymptotically as:
        \[
          \frac{r^{2-2a}}{1+r} \sim \frac{r^{2-2a}}{r} = r^{1-2a}.
        \]
        By the $p$-test for integrals, $\int_1^\infty r^p dr$ converges if and only if $p < -1$. Setting $p = 1 - 2a$:
        \[
          1 - 2a < -1 \implies 2a > 2 \implies a > 1.
        \]
        Thus, the improper integral converges for all real numbers $a > 1$.

  \item % Solution 2
        Let us represent the vectors $w_1, w_2, w_3, w_4$ as column vectors in the basis $(v_1, v_2, v_3, v_4)$. We form the coordinate matrix $A$:
        \[
          A = \begin{pmatrix} 1 & 1 & 2 & 1 \\ -1 & a-1 & -2 & -1 \\ 1 & 1 & a+2 & 1 \\ 0 & -a & 0 & a+1 \end{pmatrix}.
        \]
        \textit{(Note: Matrix representations commonly stack the coordinates as columns, but taking the transpose and using row operations gives the same rank and determinant.)} Let's take the transpose and use row operations, stacking them as rows:
        \[
          A^\top = \begin{pmatrix} 1 & -1 & 1 & 0 \\ 1 & a-1 & 1 & -a \\ 2 & -2 & a+2 & 0 \\ 1 & -1 & 1 & a+1 \end{pmatrix}.
        \]
        We apply Gaussian elimination to reduce $A^\top$ to an upper triangular form.
        Subtract Row 1 from Row 2; subtract $2 \times$ Row 1 from Row 3; subtract Row 1 from Row 4:
        \[
          \sim \begin{pmatrix} 1 & -1 & 1 & 0 \\ 0 & a & 0 & -a \\ 0 & 0 & a & 0 \\ 0 & 0 & 0 & a+1 \end{pmatrix}.
        \]
        \begin{enumerate}[label=(\arabic*)]
          \item The vectors $w_1, w_2, w_3, w_4$ are linearly independent if and only if $\det(A) = \det(A^\top) \neq 0$.
          Since the reduced matrix is upper triangular, its determinant is the product of its diagonal entries:
          \[
            \det(A^\top) = 1 \cdot a \cdot a \cdot (a+1) = a^2(a+1).
          \]
          Therefore, the vectors are linearly independent when $a^2(a+1) \neq 0$, which means $a \neq 0$ and $a \neq -1$.
          
          \item The vectors are linearly dependent when $a = 0$ or $a = -1$.
          \begin{itemize}
            \item \textbf{Case $a = 0$:} Substituting $a = 0$ into the reduced matrix yields:
            \[
              \begin{pmatrix} 1 & -1 & 1 & 0 \\ 0 & 0 & 0 & 0 \\ 0 & 0 & 0 & 0 \\ 0 & 0 & 0 & 1 \end{pmatrix}.
            \]
            There are exactly $2$ non-zero rows (Row 1 and Row 4). Thus, the rank of the matrix is $2$, meaning $\dim(\langle w_1, w_2, w_3, w_4 \rangle) = 2$.
            
            \item \textbf{Case $a = -1$:} Substituting $a = -1$ into the reduced matrix yields:
            \[
              \begin{pmatrix} 1 & -1 & 1 & 0 \\ 0 & -1 & 0 & 1 \\ 0 & 0 & -1 & 0 \\ 0 & 0 & 0 & 0 \end{pmatrix}.
            \]
            There are exactly $3$ non-zero rows (Rows 1, 2, and 3). Thus, the rank of the matrix is $3$, meaning $\dim(\langle w_1, w_2, w_3, w_4 \rangle) = 3$.
          \end{itemize}
        \end{enumerate}

  \item % Solution 3
        The group $G = \mathbb{Z}_2 \times \mathbb{Z}_3 \times \mathbb{Z}_3 \times \mathbb{Z}_4$ has order $|G| = 2 \times 3 \times 3 \times 4 = 72$.
        The subgroup $H$ is isomorphic to $\mathbb{Z}_2 \times \mathbb{Z}_3 \cong \mathbb{Z}_6$, so $|H| = 6$.
        We are given a subgroup $K$ of order $|K| = 12$ such that $G = K + H$.
        For finite groups, $|K + H| = \frac{|K||H|}{|K \cap H|}$.
        We require $K + H = G$, meaning $|K + H| = 72$. Thus,
        \[
          72 = \frac{12 \times 6}{|K \cap H|} \implies |K \cap H| = 1 \implies K \cap H = \{0\}.
        \]
        This implies $G$ is the direct sum $G = K \oplus H$, making $K$ a complement of $H$ in $G$.
        
        We can decompose $G$ into its Sylow $p$-subgroups (which are just the $p$-torsion components since $G$ is abelian):
        $G = G_2 \times G_3$, where $G_2 = \mathbb{Z}_2 \times \mathbb{Z}_4$ and $G_3 = \mathbb{Z}_3 \times \mathbb{Z}_3$.
        Correspondingly, $H = H_2 \times H_3$, where $H_2 = \mathbb{Z}_2 \times \{0\} \le G_2$ and $H_3 = \mathbb{Z}_3 \times \{0\} \le G_3$.
        Any subgroup $K \le G$ of order $12$ decomposes as $K = K_2 \times K_3$ with $K_2 \le G_2$ and $K_3 \le G_3$.
        The direct sum condition $G = K \oplus H$ decouples into $G_2 = K_2 \oplus H_2$ and $G_3 = K_3 \oplus H_3$.
        
        \textbf{Finding $K_2$ in $G_2$:}
        We need a complement $K_2$ to $H_2 = \langle(1, 0)\rangle$ in $\mathbb{Z}_2 \times \mathbb{Z}_4$.
        Since $|G_2| = 8$ and $|H_2| = 2$, we need $|K_2| = 4$. Furthermore, $G_2 / H_2 \cong \mathbb{Z}_4$, so $K_2$ must be isomorphic to $\mathbb{Z}_4$.
        Thus, $K_2$ is generated by an element of order $4$, which must be of the form $(x, y)$ where $y \in \{1, 3\}$.
        If $y = 1$, the generator is either $(0, 1)$ or $(1, 1)$. 
        \begin{itemize}
          \item $\langle (0, 1) \rangle = \{(0,0), (0,1), (0,2), (0,3)\}$. Its intersection with $H_2$ is $\{(0,0)\}$.
          \item $\langle (1, 1) \rangle = \{(0,0), (1,1), (0,2), (1,3)\}$. Its intersection with $H_2$ is $\{(0,0)\}$.
        \end{itemize}
        The generators $(0,3)$ and $(1,3)$ generate the exact same two subgroups. Thus, there are exactly $2$ valid choices for $K_2$.

        \textbf{Finding $K_3$ in $G_3$:}
        We need a complement $K_3$ to $H_3 = \langle(1, 0)\rangle$ in $\mathbb{Z}_3 \times \mathbb{Z}_3$.
        This is a vector space of dimension $2$ over the field $\mathbb{F}_3$. $H_3$ is a $1$-dimensional subspace, and we need another $1$-dimensional subspace $K_3$ such that they span the whole space.
        The total number of $1$-dimensional subspaces in $\mathbb{F}_3^2$ is $\frac{3^2 - 1}{3 - 1} = 4$.
        One of them is $H_3$. The remaining $3$ subspaces are valid complements (they are generated by $(0,1), (1,1)$, and $(2,1)$). Thus, there are exactly $3$ valid choices for $K_3$.

        The total number of such subgroups $K$ is the product of the number of choices for $K_2$ and $K_3$, which is $2 \times 3 = 6$.

  \item % Solution 4
        Using the reverse triangle inequality $|a + b| \ge |a| - |b|$ and the fact that $|\sin t| \le 1$, we can bound the integrand from below for any $x \in \mathbb{R}$ and $t \ge 0$:
        \[
          |x + t \sin t| e^{-t} \ge (|x| - |t \sin t|) e^{-t} \ge (|x| - t) e^{-t}.
        \]
        Integrating this lower bound over $t \in [0, \infty)$ gives:
        \[
          f(x) \ge \int_0^\infty (|x| - t) e^{-t} dt = |x| \int_0^\infty e^{-t} dt - \int_0^\infty t e^{-t} dt.
        \]
        Using the standard integral evaluations $\int_0^\infty e^{-t} dt = 1$ and $\int_0^\infty t e^{-t} dt = 1$ (via integration by parts), we get:
        \[
          f(x) \ge |x| - 1.
        \]
        From this, it is evident that as $|x| \to \infty$, $f(x) \to \infty$. Specifically, for $|x| \ge f(0) + 1$, we are guaranteed that $f(x) \ge f(0)$. This means that the minimum of $f$ on $\mathbb{R}$, if it exists, must occur within the compact interval $[-f(0)-1, f(0)+1]$.

        To ensure a minimum is attained on this compact interval, we must show that $f(x)$ is continuous. We evaluate the difference $|f(x) - f(y)|$:
        \begin{align*}
          |f(x) - f(y)| &= \left| \int_0^\infty (|x + t \sin t| - |y + t \sin t|) e^{-t} dt \right| \\
          &\le \int_0^\infty \big| |x + t \sin t| - |y + t \sin t| \big| e^{-t} dt \\
          &\le \int_0^\infty |(x + t \sin t) - (y + t \sin t)| e^{-t} dt \\
          &= \int_0^\infty |x - y| e^{-t} dt = |x - y|.
        \end{align*}
        Thus, $f(x)$ is globally Lipschitz continuous (with Lipschitz constant $1$), which strongly implies it is continuous. 
        Since $f(x)$ is continuous on $\mathbb{R}$ and coercive ($\lim_{|x| \to \infty} f(x) = \infty$), by the Extreme Value Theorem, $f(x)$ must achieve a global minimum value on $\mathbb{R}$.

  \item % Solution 5
        To find the value of the improper integral $I = \int_{-\infty}^\infty \frac{x^3 \sin x}{(x^2 + 1)^2} dx$, we will evaluate the associated contour integral
        \[
          \oint_C g(z) dz, \quad \text{where } g(z) = \frac{z^3 e^{iz}}{(z^2 + 1)^2}
        \]
        over the standard semicircular contour $C$ in the upper half-plane. The contour consists of the real segment $[-R, R]$ and the semicircular arc $\Gamma_R : z = R e^{i\theta}, 0 \le \theta \le \pi$.
        The function $g(z)$ has poles where $(z^2 + 1)^2 = 0$, which are poles of order 2 at $z = i$ and $z = -i$. Only $z = i$ lies inside our contour $C$ (in the upper half-plane).
        By Jordan's Lemma, the integral over the arc $\Gamma_R$ vanishes as $R \to \infty$ because $\frac{z^3}{(z^2+1)^2} \to 0$ uniformly as $|z| \to \infty$. Thus:
        \[
          \int_{-\infty}^\infty \frac{x^3 e^{ix}}{(x^2 + 1)^2} dx = 2\pi i \operatorname{Res}(g(z), i).
        \]
        To compute the residue at the pole of order 2 at $z = i$:
        \[
          \operatorname{Res}(g(z), i) = \lim_{z \to i} \frac{d}{dz} \left[ (z - i)^2 g(z) \right] = \lim_{z \to i} \frac{d}{dz} \left[ \frac{z^3 e^{iz}}{(z + i)^2} \right].
        \]
        Let $h(z) = \frac{z^3 e^{iz}}{(z + i)^2}$. Using logarithmic differentiation $\frac{h'(z)}{h(z)} = \frac{3}{z} + i - \frac{2}{z+i}$, we find:
        \[
          h'(i) = h(i) \left( \frac{3}{i} + i - \frac{2}{i+i} \right).
        \]
        Evaluating $h(i)$ and the terms inside the parenthesis:
        \begin{align*}
          h(i) &= \frac{i^3 e^{i(i)}}{(i+i)^2} = \frac{-i e^{-1}}{(2i)^2} = \frac{-i e^{-1}}{-4} = \frac{i}{4e}, \\
          \left( \frac{3}{i} + i - \frac{2}{2i} \right) &= -3i + i + i = -i.
        \end{align*}
        Thus, the residue is:
        \[
          \operatorname{Res}(g(z), i) = h'(i) = \left( \frac{i}{4e} \right) (-i) = \frac{1}{4e}.
        \]
        The value of the contour integral is therefore:
        \[
          \int_{-\infty}^\infty \frac{x^3 e^{ix}}{(x^2 + 1)^2} dx = 2\pi i \left( \frac{1}{4e} \right) = \frac{\pi i}{2e}.
        \]
        Using Euler's formula $e^{ix} = \cos x + i \sin x$, we take the imaginary part of both sides:
        \[
          \int_{-\infty}^\infty \frac{x^3 \sin x}{(x^2 + 1)^2} dx = \operatorname{Im}\left( \frac{\pi i}{2e} \right) = \frac{\pi}{2e}.
        \]

  \item % Solution 6
        Let $F : \mathbb{R}^4 \to \mathbb{R}^2$ be the map defined by
        \[
          F(x, y, z, w) = (x^2 + y^2 + z^2 + w^2 - 1, x^3 + y^3 + z^3 + w^3).
        \]
        The subspace $X$ is exactly the level set $F^{-1}(0, 0)$. 
        First, observe that $X$ is a closed subset of the unit sphere $S^3 \subset \mathbb{R}^4$ (since it is the intersection of $S^3$ with the closed set $x^3 + y^3 + z^3 + w^3 = 0$). Since $S^3$ is compact (closed and bounded by the Heine-Borel theorem), its closed subset $X$ must also be compact.
        
        To prove that $X$ is a differentiable manifold, by the Regular Level Set Theorem, it suffices to show that $(0, 0)$ is a regular value of $F$. This means the Jacobian matrix $dF$ must have full rank (rank 2) at every point in $X$.
        The Jacobian matrix $dF$ is given by:
        \[
          dF = \begin{pmatrix} 2x & 2y & 2z & 2w \\ 3x^2 & 3y^2 & 3z^2 & 3w^2 \end{pmatrix}.
        \]
        Assume for the sake of contradiction that the rank of $dF$ is strictly less than 2 at some point $P = (x, y, z, w) \in X$. This implies the two rows are linearly dependent. Thus, there exists a scalar $\lambda \in \mathbb{R}$ such that:
        \[
          (3x^2, 3y^2, 3z^2, 3w^2) = \lambda (2x, 2y, 2z, 2w).
        \]
        For each variable $v \in \{x, y, z, w\}$, we have $3v^2 = 2\lambda v$, which implies either $v = 0$ or $v = \frac{2\lambda}{3}$. Let $c = \frac{2\lambda}{3}$. Thus, every coordinate is restricted to the set $\{0, c\}$.
        Let $k$ be the number of coordinates that are equal to $c$ (so $4 - k$ coordinates are $0$).
        Since $P \in X$, it must satisfy $x^2 + y^2 + z^2 + w^2 = 1$. This means we cannot have all coordinates equal to $0$. Hence, $k \ge 1$ and $c \neq 0$.
        $P$ must also satisfy the second equation $x^3 + y^3 + z^3 + w^3 = 0$. Substituting our coordinates yields:
        \[
          k \cdot c^3 = 0.
        \]
        Since $k \ge 1$, we are forced to conclude $c^3 = 0$, giving $c = 0$. This contradicts our earlier deduction that $c \neq 0$.
        Therefore, our assumption that the rows are linearly dependent must be false. The Jacobian matrix $dF$ has full rank everywhere on $X$.
        Consequently, $X$ is a regular submanifold of $\mathbb{R}^4$ of dimension $4 - 2 = 2$, proving it is a compact differentiable manifold.
\end{enumerate}

\end{document}
